\documentclass{article}
\usepackage[utf8]{inputenc}
\usepackage[T1]{fontenc}
\usepackage[english, polish]{babel}
\usepackage{amsmath}
\usepackage{amsfonts}
\usepackage{amssymb}
\usepackage{graphicx}
\usepackage{float}
\usepackage{parskip}
\usepackage{hyperref}
% Hide the link borders
\hypersetup{hidelinks}
\usepackage{geometry}
% Margins for A4 paper
\geometry{a4paper, margin=2cm}
\usepackage{enumitem}
% Spacing for itemize lists
\setlist[itemize]{itemsep=1em, topsep=0.25em, parsep=0pt, partopsep=0pt}
\setlength{\parindent}{0pt} % No paragraph indentation
\setlength{\emergencystretch}{2em} % Prevent overfull hboxes

\title{Sprawozdanie z zadania 6\\\textbf{Rozszerzenie gry wyścigowej}}
\author{Mikołaj Kuryło s193743\\Antonina Włodarczyk s197793}
\date{}

\begin{document}
\maketitle
\vspace{1cm}

\section{Sformułowanie zadania i przyjęte założenia}
Celem projektu było rozbudowanie świata gry poprzez dodanie możliwości zmiany generowanych map.
Mapy te umożliwiają ciekawszą rozgrywkę i dzięki uwzględnieniu

\section{Opis modelu map}
Mapy zostają dodane w programie dzięki plikom o rozszerzeniu json.
Jest to możliwe, dzięki pliku map\_reader, który służy jako translator. 

\subsection{Przykładowy kod jednej z map}
\begin{verbatim}
   {
  "name": "map_3",
  "width": 5,
  "height": 9,
  "tile_friction": {
    "_": 0.05,
    "X": 0.9
  },
  "start_position": [4, 1],
  "end_position": [4, 8],
  "map": [
    ["_", "X", "X", "X", "X"],
    ["X", "X", "_", "X", "X"],
    ["X", "X", "_", "_", "_"],
    ["X", "X", "X", "X", "_"],
    ["_", "_", "_", "X", "_"],
    ["_", "X", "X", "X", "_"],
    ["_", "X", "_", "_", "_"],
    ["_", "X", "X", "X", "X"],
    ["_", "_", "_", "_", "_"]
  ]
}

\end{verbatim}
Mapy z tej postaci są czytane głównie dzięki kodu odpowiedzialnemu za konwersję.
\begin{verbatim}
class MapReader():
    def read_map(self, file_path: str,
                 file_type: MapFileType) -> MapData:

        match file_type:
            case MapFileType.JSON:
                return self._read_json_map(file_path)
            case _:
                raise ValueError(f"Unsupported map file type: {file_type}")

    def _read_json_map(self, file_path: str) -> MapData:
        try:
            with open(file_path, 'r') as f:
                data = json.load(f)

                # Store symbolic map as-is
                symbolic_map_matrix = np.matrix(data["map"], dtype=object)

                # Create traction map by replacing symbols with friction coefficients
                tile_friction = data["tile_friction"]
                traction_grid = [
                    [tile_friction[symbol] for symbol in row]
                    for row in data["map"]
                ]
                traction_map_matrix = np.matrix(traction_grid, dtype=float)

                # Extract start and end positions
                start_position = Coords(
                    data["start_position"][0], data["start_position"][1])
                end_position = Coords(
                    data["end_position"][0], data["end_position"][1])

                return MapData(
                    symbolic_map_matrix=symbolic_map_matrix,
                    traction_map_matrix=traction_map_matrix,
                    start_position=start_position,
                    end_position=end_position,
                    size_x=int(len(data["map"][0]) * const.MAP_SCALE_FACTOR),
                    size_y=int(len(data["map"]) * const.MAP_SCALE_FACTOR)
                )

        except FileNotFoundError:
            print(f"Error: Map file '{file_path}' not found.")
            raise
        except KeyError as e:
            print(f"Error: Missing key in map file: {e}")
            raise
        except json.JSONDecodeError:
            print(f"Error: Map file '{file_path}' is not valid JSON.")
            raise
        except Exception as e:
            print(
                f"Error: Unexpected error reading map file '{file_path}': {e}")
            raise

\end{verbatim}

\section{Dyskusja osiągniętych wyników}
\subsection{Zalety}
Wykorzystywanie pliku map\_reader znacznie ułatwia i uproszcza dodawanie rozbudowanych oraz znacznie dłuższych tras.
Jest to prosty proces, który nie wymaga wielu modyfikacji.

\subsection{Ograniczenia i potencjalne ulepszenia}
Mapy nadal pozostają nieruchome, oraz nie zawsze jest poprowadzona droga, która wskazuje najszybszą trasę. W niektórych przypadkach może to wprowadzać gracza w zdziwienie, ponieważ cel nie zawsze jest okreslony. W przypadku gdzie trasa znajduje sie w centrum planszy, a na zewnątrz znajdują się elementy tła, to cel może być gdziekolwiek.
Nie ułatwia to gry, a wyjechanie w nieodpowiednim miejsu może skutkowac zakończeniem gry z powodu wypadnięcia za mapę.

\end{document}